
\documentclass{article}
\usepackage{colortbl}
\usepackage{makecell}
\usepackage{multirow}
\usepackage{supertabular}

\begin{document}

\newcounter{utterance}

\centering \large Interaction Transcript for game `cladder', experiment `full\_v1.5\_default', episode 8797 with refineonly{-}Qwen3.5{-}27B.
\vspace{24pt}

{ \footnotesize  \setcounter{utterance}{1}
\setlength{\tabcolsep}{0pt}
\begin{supertabular}{c@{$\;$}|p{.15\linewidth}@{}p{.15\linewidth}p{.15\linewidth}p{.15\linewidth}p{.15\linewidth}p{.15\linewidth}}
    \# & \multicolumn{2}{c}{Player} && \multicolumn{2}{c}{Game Master} \\
    \hline

    \theutterance \stepcounter{utterance}  
    & & & \multicolumn{4}{p{0.6\linewidth}}{
        \cellcolor[rgb]{0.9,0.9,0.9}{
            \makecell[{{p{\linewidth}}}]{
                \texttt{\tiny{[P1$\langle$GM]}}
                \texttt{You are an expert in causal inference. The following question is not a typical commonsense query, but rather a meticulously designed question created by a professor specializing in causal inference, intended to assess the students' mastery of the course content.} \\
\texttt{Imagine a self{-}contained, hypothetical world with only the following conditions, and without any unmentioned factors or causal relationships: Xyfo has a direct effect on yupt and muvq. Yupt has a direct effect on zupj. Zupj has a direct effect on muvq. Xyfo is unobserved. For those who are not yupt and are not zupj, the probability of muvq is 90\%. For those who are not yupt and are zupj, the probability of muvq is 54\%. For those who are yupt and are not zupj, the probability of muvq is 96\%. For those who are yupt and are zupj, the probability of muvq is 57\%. For those who are not yupt, the probability of zupj is 51\%. For those who are yupt, the probability of zupj is 34\%. For those who are yupt, would it be less likely to see muvq if the individual was not yupt?} \\
\texttt{Start your answer with 'Yes' or 'No', followed by additional reasoning or evidence to support your explanation.} \\
            }
        }
    }
     \\ \\

    \theutterance \stepcounter{utterance}  
    & & \multicolumn{4}{p{0.6\linewidth}}{
        \cellcolor[rgb]{0.95,0.95,0.95}{
            \makecell[{{p{\linewidth}}}]{
                \texttt{\tiny{[P1$\leftrightarrow$P]}}
                \texttt{{[}DRAFT{]}} \\
\texttt{No} \\
\\ 
\texttt{**Reasoning:**} \\
\\ 
\texttt{To answer this question, we must calculate the **Average Treatment Effect (ATE)** or, more specifically in this context, the counterfactual probability of the outcome (\$muvq\$) for the group defined as "those who are \$yupt\$" under the hypothetical intervention where they are "not \$yupt\$".} \\
\\ 
\texttt{Let \$Y\$ represent \$yupt\$, \$Z\$ represent \$zupj\$, and \$M\$ represent \$muvq\$.} \\
\texttt{The question asks: For individuals where \$Y=1\$, is \$P(M=1 \textbackslash{}mid do(Y=0)) < P(M=1 \textbackslash{}mid Y=1)\$?} \\
\\ 
\texttt{**Step 1: Calculate the observed probability of \$muvq\$ for those who are \$yupt\$ (\$P(M=1 \textbackslash{}mid Y=1)\$).**} \\
\texttt{We use the law of total probability, conditioning on \$Z\$:} \\
\texttt{\$\$P(M=1 \textbackslash{}mid Y=1) = P(M=1 \textbackslash{}mid Y=1, Z=1)P(Z=1 \textbackslash{}mid Y=1) + P(M=1 \textbackslash{}mid Y=1, Z=0)P(Z=0 \textbackslash{}mid Y=1)\$\$} \\
\\ 
\texttt{From the prompt:} \\
\texttt{*   \$P(M=1 \textbackslash{}mid Y=1, Z=1) = 0.57\$} \\
\texttt{*   \$P(M=1 \textbackslash{}mid Y=1, Z=0) = 0.96\$} \\
\texttt{*   \$P(Z=1 \textbackslash{}mid Y=1) = 0.34\$} \\
\texttt{*   Therefore, \$P(Z=0 \textbackslash{}mid Y=1) = 1 {-} 0.34 = 0.66\$} \\
\\ 
\texttt{Calculation:} \\
\texttt{\$\$P(M=1 \textbackslash{}mid Y=1) = (0.57 \textbackslash{}times 0.34) + (0.96 \textbackslash{}times 0.66)\$\$} \\
\texttt{\$\$P(M=1 \textbackslash{}mid Y=1) = 0.1938 + 0.6336 = 0.8274\$\$} \\
\\ 
\texttt{**Step 2: Calculate the counterfactual probability of \$muvq\$ for those who are \$yupt\$ if they were not \$yupt\$ (\$P(M=1 \textbackslash{}mid do(Y=0))\$).**} \\
\texttt{In causal inference, \$P(M=1 \textbackslash{}mid do(Y=0))\$ represents the probability of the outcome if we intervene to set \$Y=0\$ for the entire population (or the specific subgroup, assuming the structural equations remain consistent). Since \$Xyfo\$ is unobserved but acts as a common cause, we must be careful. However, the prompt states "without any unmentioned factors or causal relationships" and gives us the conditional probabilities of \$Z\$ given \$Y\$ and \$M\$ given \$Y,} \\
            }
        }
    }
    & \\ \\

    \theutterance \stepcounter{utterance}  
    & \multicolumn{4}{p{0.6\linewidth}}{
        \cellcolor[rgb]{0.561,0.737,0.561}{
            \makecell[{{p{\linewidth}}}]{
                \texttt{\tiny{[P1$\rangle$GM]}}
                \texttt{No} \\
\\ 
\texttt{**Reasoning:**} \\
\\ 
\texttt{To answer this question, we must compare the observed probability of \$muvq\$ for individuals who are \$yupt\$ against the counterfactual probability of \$muvq\$ for those same individuals if they were not \$yupt\$.} \\
\\ 
\texttt{Let \$Y\$ represent \$yupt\$, \$Z\$ represent \$zupj\$, and \$M\$ represent \$muvq\$.} \\
\texttt{The question asks: For individuals where \$Y=1\$, is \$P(M=1 \textbackslash{}mid do(Y=0)) < P(M=1 \textbackslash{}mid Y=1)\$?} \\
\\ 
\texttt{**Step 1: Calculate the observed probability of \$muvq\$ for those who are \$yupt\$ (\$P(M=1 \textbackslash{}mid Y=1)\$).**} \\
\texttt{We use the law of total probability, conditioning on \$Z\$:} \\
\texttt{\$\$P(M=1 \textbackslash{}mid Y=1) = P(M=1 \textbackslash{}mid Y=1, Z=1)P(Z=1 \textbackslash{}mid Y=1) + P(M=1 \textbackslash{}mid Y=1, Z=0)P(Z=0 \textbackslash{}mid Y=1)\$\$} \\
\\ 
\texttt{From the prompt:} \\
\texttt{*   \$P(M=1 \textbackslash{}mid Y=1, Z=1) = 0.57\$} \\
\texttt{*   \$P(M=1 \textbackslash{}mid Y=1, Z=0) = 0.96\$} \\
\texttt{*   \$P(Z=1 \textbackslash{}mid Y=1) = 0.34\$} \\
\texttt{*   Therefore, \$P(Z=0 \textbackslash{}mid Y=1) = 1 {-} 0.34 = 0.66\$} \\
\\ 
\texttt{Calculation:} \\
\texttt{\$\$P(M=1 \textbackslash{}mid Y=1) = (0.57 \textbackslash{}times 0.34) + (0.96 \textbackslash{}times 0.66)\$\$} \\
\texttt{\$\$P(M=1 \textbackslash{}mid Y=1) = 0.1938 + 0.6336 = 0.8274\$\$} \\
\\ 
\texttt{**Step 2: Calculate the counterfactual probability of \$muvq\$ for those who are \$yupt\$ if they were not \$yupt\$ (\$P(M=1 \textbackslash{}mid do(Y=0))\$).**} \\
\texttt{The prompt specifies a self{-}contained world with no unmentioned factors. The causal structure is \$Xyfo \textbackslash{}to Y\$, \$Xyfo \textbackslash{}to M\$, \$Y \textbackslash{}to Z\$, \$Z \textbackslash{}to M\$, and \$Y \textbackslash{}to M\$ (implied by the direct effect mentioned).} \\
\texttt{When we intervene to set \$Y=0\$ (\$do(Y=0)\$), we sever the incoming arrow from \$Xyfo\$ to \$Y\$. However, the distribution of \$Z\$ and \$M\$ given \$Y=0\$ is determined by the structural equations governing the system when \$Y\$} \\
            }
        }
    }
    & & \\ \\

    \theutterance \stepcounter{utterance}  
    & & & \multicolumn{2}{p{0.3\linewidth}}{
        \cellcolor[rgb]{0.9,0.9,0.9}{
            \makecell[{{p{\linewidth}}}]{
                \texttt{\tiny{[GM$|$GM]}}
                \texttt{No} \\
            }
        }
    }
    & & \\ \\

    \theutterance \stepcounter{utterance}  
    & & & \multicolumn{2}{p{0.3\linewidth}}{
        \cellcolor[rgb]{0.9,0.9,0.9}{
            \makecell[{{p{\linewidth}}}]{
                \texttt{\tiny{[GM$|$GM]}}
                \texttt{yes} \\
            }
        }
    }
    & & \\ \\

    \theutterance \stepcounter{utterance}  
    & & & \multicolumn{2}{p{0.3\linewidth}}{
        \cellcolor[rgb]{0.9,0.9,0.9}{
            \makecell[{{p{\linewidth}}}]{
                \texttt{\tiny{[GM$|$GM]}}
                \texttt{game\_result = LOSE} \\
            }
        }
    }
    & & \\ \\

\end{supertabular}
}

\end{document}
